\section{Tutorial 1: Foundations}

\subsection{Objectives}

This tutorial reinforces the key ideas from Lectures 1 and 2:
\begin{itemize}
    \item Learning as function approximation
    \item Probabilistic formulation of learning
    \item Empirical risk minimization (ERM)
    \item Bias--variance tradeoff (intuition)
\end{itemize}

\medskip

\noindent
\textbf{Focus.}  
Develop intuition through concrete examples and simple derivations.

\subsection{Exercise 1: Empirical Risk vs True Risk}

Let $(x,y) \sim P$, and let $f$ be a fixed model.

\medskip

\noindent
\textbf{Tasks.}
\begin{enumerate}
    \item Define the true risk $R(f)$.
    \item Define the empirical risk $\hat{R}_n(f)$.
    \item Explain why $\hat{R}_n(f)$ is a good approximation of $R(f)$.
    \item Give an example where $\hat{R}_n(f)$ is misleading.
\end{enumerate}

\medskip

\noindent
\textbf{Solution.}

\begin{enumerate}
    \item 
    \[
    R(f) = \mathbb{E}_{(x,y)\sim P}[\ell(f(x),y)].
    \]

    \item 
    \[
    \hat{R}_n(f) = \frac{1}{n} \sum_{i=1}^n \ell(f(x_i), y_i).
    \]

    \item By the law of large numbers:
    \[
    \hat{R}_n(f) \to R(f) \quad \text{as } n \to \infty.
    \]

    \item If $f$ overfits the training data (e.g., memorizes noise), then:
    \[
    \hat{R}_n(f) \ll R(f).
    \]
\end{enumerate}

\medskip

\noindent
\textbf{Key takeaway.}  
Empirical risk approximates true risk, but only when the model generalizes.

\subsection{Exercise 2: Bias--Variance Intuition}

Consider fitting a model to noisy data:
\[
y = f(x) + \varepsilon, \quad \varepsilon \sim \mathcal{N}(0, \sigma^2).
\]

\medskip

\noindent
\textbf{Tasks.}
\begin{enumerate}
    \item Describe what happens when the model is too simple.
    \item Describe what happens when the model is too complex.
    \item Explain the bias--variance tradeoff qualitatively.
\end{enumerate}

\medskip

\noindent
\textbf{Solution.}

\begin{enumerate}
    \item \textbf{Too simple (high bias):}
    \begin{itemize}
        \item Model cannot capture true structure
        \item Systematic error remains
    \end{itemize}

    \item \textbf{Too complex (high variance):}
    \begin{itemize}
        \item Model fits noise
        \item Predictions vary significantly with data
    \end{itemize}

    \item \textbf{Tradeoff:}
    \begin{itemize}
        \item Increasing complexity reduces bias
        \item Increasing complexity increases variance
    \end{itemize}
\end{enumerate}

\medskip

\noindent
\textbf{Key takeaway.}  
Good models balance approximation accuracy and stability.

\subsection{Exercise 3: Loss Functions and Probabilistic Interpretation}

Suppose we assume:
\[
y = f_\theta(x) + \varepsilon, \quad \varepsilon \sim \mathcal{N}(0, \sigma^2).
\]

\medskip

\noindent
\textbf{Tasks.}
\begin{enumerate}
    \item Write down the likelihood $p(y \mid x, \theta)$.
    \item Show that maximizing likelihood is equivalent to minimizing squared loss.
\end{enumerate}

\medskip

\noindent
\textbf{Solution.}

\begin{enumerate}
    \item 
    \[
    p(y \mid x, \theta)
    =
    \frac{1}{\sqrt{2\pi \sigma^2}}
    \exp\left(-\frac{(y - f_\theta(x))^2}{2\sigma^2}\right).
    \]

    \item Taking negative log-likelihood:
    \[
    -\log p(y \mid x, \theta)
    =
    \frac{(y - f_\theta(x))^2}{2\sigma^2} + C.
    \]

    Minimizing this is equivalent to minimizing:
    \[
    (y - f_\theta(x))^2.
    \]
\end{enumerate}

\medskip

\noindent
\textbf{Key takeaway.}  
Loss functions encode assumptions about the data-generating process.

\subsection{Exercise 4: Simple Optimization Problem}

Consider linear regression:
\[
L(w) = \sum_{i=1}^n (w^\top x_i - y_i)^2.
\]

\medskip

\noindent
\textbf{Tasks.}
\begin{enumerate}
    \item Compute the gradient $\nabla_w L(w)$.
    \item Write the gradient descent update rule.
\end{enumerate}

\medskip

\noindent
\textbf{Solution.}

\begin{enumerate}
    \item 
    \[
    \nabla_w L(w) = 2 \sum_{i=1}^n (w^\top x_i - y_i) x_i.
    \]

    In matrix form:
    \[
    \nabla_w L(w) = 2 X^\top (Xw - y).
    \]

    \item Gradient descent update:
    \[
    w_{t+1} = w_t - \eta \nabla_w L(w_t).
    \]
\end{enumerate}

\medskip

\noindent
\textbf{Key takeaway.}  
Learning reduces to solving an optimization problem.

\subsection{Exercise 5 (Discussion): What Enables Generalization?}

\textbf{Question.}  
If many functions fit the data perfectly, why do some generalize better than others?

\medskip

\noindent
\textbf{Discussion points.}
\begin{itemize}
    \item Hypothesis space restrictions
    \item Inductive bias
    \item Simplicity of solutions
\end{itemize}

\medskip

\noindent
\textbf{Insight.}  
Generalization depends on assumptions beyond data.

\subsection{Summary}

\begin{itemize}
    \item ERM approximates true risk but may fail under overfitting
    \item Bias--variance tradeoff governs model complexity
    \item Loss functions arise from probabilistic assumptions
    \item Learning involves optimization over parameters
\end{itemize}

\medskip

\noindent
\textbf{Preparation for next lecture.}  
We now examine how hypothesis spaces and inductive bias formalize these ideas.